\documentclass[a4paper,12pt,oneside,oldfontcommands]{memoir}

% Castellano
% es-ucroman: evitar el warning por Fake sc para números romanos en minúscula
\usepackage[spanish,es-tabla,es-ucroman]{babel}
\selectlanguage{spanish}
\usepackage[utf8]{inputenc}
\usepackage[T1]{fontenc}
\usepackage{lmodern} % Scalable font
\usepackage{microtype}
\usepackage{placeins}

\RequirePackage{booktabs}
\RequirePackage[table]{xcolor}
\RequirePackage{xtab}
\RequirePackage{multirow}

% Links
\usepackage[colorlinks]{hyperref}
\hypersetup{
	allcolors = {blue}
}

% Ecuaciones y simbolos (tanto AMS como LaTeX
\usepackage{amsmath,amssymb,latexsym}

% Rutas de fichero / paquete
\newcommand{\ruta}[1]{{\sffamily #1}}

% Párrafos
\nonzeroparskip

% Huérfanas y viudas
\widowpenalty100000
\clubpenalty100000

% Paper size and margins (para A4)
\settrimmedsize{297mm}{210mm}{*}
\setlength{\trimtop}{0pt}
\setlength{\trimedge}{\stockwidth}
\addtolength{\trimedge}{-\paperwidth}
\settypeblocksize{634pt}{448.13pt}{*}
\setulmargins{4cm}{*}{*}
\setlrmarginsandblock{2.5cm}{2.5cm}{*}
\setmarginnotes{36pt}{36pt}{\onelineskip}
\setheadfoot{\onelineskip}{2\onelineskip}
\setheaderspaces{*}{2\onelineskip}{*}
\setlength{\headheight}{15.5pt}  % Aumenta la altura del encabezado
\checkandfixthelayout

% Colores locales
\usepackage{xcolor}
\definecolor{ctitbox}{HTML}{E6E6FF}
\definecolor{capnum}{HTML}{0606FF}

% Imagenes
\usepackage{graphicx}
\newcommand{\imagen}[2]{
	\begin{figure}[!h]
		\centering
		\includegraphics[width=0.9\textwidth]{#1}
		\caption{#2}\label{fig:#1}
	\end{figure}
	\FloatBarrier
}

\newcommand{\imagenflotante}[2]{
	\begin{figure}%[!h]
		\centering
		\includegraphics[width=0.9\textwidth]{#1}
		\caption{#2}\label{fig:#1}
	\end{figure}
}


% El comando \figura nos permite insertar figuras comodamente, y utilizando
% siempre el mismo formato. Los parametros son:
% 1 -> Porcentaje del ancho de página que ocupará la figura (de 0 a 1)
% 2 --> Fichero de la imagen
% 3 --> Texto a pie de imagen
% 4 --> Etiqueta (label) para referencias
% 5 --> Opciones que queramos pasarle al \includegraphics
% 6 --> Opciones de posicionamiento a pasarle a \begin{figure}
\newcommand{\figuraConPosicion}[6]{%
  \setlength{\anchoFloat}{#1\textwidth}%
  \addtolength{\anchoFloat}{-4\fboxsep}%
  \setlength{\anchoFigura}{\anchoFloat}%
  \begin{figure}[#6]
    \begin{center}%
      \Ovalbox{%
        \begin{minipage}{\anchoFloat}%
          \begin{center}%
            \includegraphics[width=\anchoFigura,#5]{#2}%
            \caption{#3}%
            \label{#4}%
          \end{center}%
        \end{minipage}
      }%
    \end{center}%
  \end{figure}%
}

%
% Comando para incluir imágenes en formato apaisado (sin marco).
\newcommand{\figuraApaisadaSinMarco}[5]{%
  \begin{figure}%
    \begin{center}%
    \includegraphics[angle=90,height=#1\textheight,#5]{#2}%
    \caption{#3}%
    \label{#4}%
    \end{center}%
  \end{figure}%
}
% Para las tablas
\newcommand{\otoprule}{\midrule [\heavyrulewidth]}
%
% Nuevo comando para tablas pequeñas (menos de una página).
\newcommand{\tablaSmall}[5]{%
 \begin{table}
  \begin{center}
   \rowcolors {2}{gray!35}{}
   \begin{tabular}{#2}
    \toprule
    #4
    \otoprule
    #5
    \bottomrule
   \end{tabular}
   \caption{#1}
   \label{tabla:#3}
  \end{center}
 \end{table}
}

%
% Nuevo comando para tablas pequeñas (menos de una página).
\newcommand{\tablaSmallSinColores}[5]{%
 \begin{table}[H]
  \begin{center}
   \begin{tabular}{#2}
    \toprule
    #4
    \otoprule
    #5
    \bottomrule
   \end{tabular}
   \caption{#1}
   \label{tabla:#3}
  \end{center}
 \end{table}
}

\newcommand{\tablaApaisadaSmall}[5]{%
\begin{landscape}
  \begin{table}
   \begin{center}
    \rowcolors {2}{gray!35}{}
    \begin{tabular}{#2}
     \toprule
     #4
     \otoprule
     #5
     \bottomrule
    \end{tabular}
    \caption{#1}
    \label{tabla:#3}
   \end{center}
  \end{table}
\end{landscape}
}

%
% Nuevo comando para tablas grandes con cabecera y filas alternas coloreadas en gris.
\newcommand{\tabla}[6]{%
  \begin{center}
    \tablefirsthead{
      \toprule
      #5
      \otoprule
    }
    \tablehead{
      \multicolumn{#3}{l}{\small\sl continúa desde la página anterior}\\
      \toprule
      #5
      \otoprule
    }
    \tabletail{
      \hline
      \multicolumn{#3}{r}{\small\sl continúa en la página siguiente}\\
    }
    \tablelasttail{
      \hline
    }
    \bottomcaption{#1}
    \rowcolors {2}{gray!35}{}
    \begin{xtabular}{#2}
      #6
      \bottomrule
    \end{xtabular}
    \label{tabla:#4}
  \end{center}
}

%
% Nuevo comando para tablas grandes con cabecera.
\newcommand{\tablaSinColores}[6]{%
  \begin{center}
    \tablefirsthead{
      \toprule
      #5
      \otoprule
    }
    \tablehead{
      \multicolumn{#3}{l}{\small\sl continúa desde la página anterior}\\
      \toprule
      #5
      \otoprule
    }
    \tabletail{
      \hline
      \multicolumn{#3}{r}{\small\sl continúa en la página siguiente}\\
    }
    \tablelasttail{
      \hline
    }
    \bottomcaption{#1}
    \begin{xtabular}{#2}
      #6
      \bottomrule
    \end{xtabular}
    \label{tabla:#4}
  \end{center}
}

%
% Nuevo comando para tablas grandes sin cabecera.
\newcommand{\tablaSinCabecera}[5]{%
  \begin{center}
    \tablefirsthead{
      \toprule
    }
    \tablehead{
      \multicolumn{#3}{l}{\small\sl continúa desde la página anterior}\\
      \hline
    }
    \tabletail{
      \hline
      \multicolumn{#3}{r}{\small\sl continúa en la página siguiente}\\
    }
    \tablelasttail{
      \hline
    }
    \bottomcaption{#1}
  \begin{xtabular}{#2}
    #5
   \bottomrule
  \end{xtabular}
  \label{tabla:#4}
  \end{center}
}



\definecolor{cgoLight}{HTML}{EEEEEE}
\definecolor{cgoExtralight}{HTML}{FFFFFF}

%
% Nuevo comando para tablas grandes sin cabecera.
\newcommand{\tablaSinCabeceraConBandas}[5]{%
  \begin{center}
    \tablefirsthead{
      \toprule
    }
    \tablehead{
      \multicolumn{#3}{l}{\small\sl continúa desde la página anterior}\\
      \hline
    }
    \tabletail{
      \hline
      \multicolumn{#3}{r}{\small\sl continúa en la página siguiente}\\
    }
    \tablelasttail{
      \hline
    }
    \bottomcaption{#1}
    \rowcolors[]{1}{cgoExtralight}{cgoLight}

  \begin{xtabular}{#2}
    #5
   \bottomrule
  \end{xtabular}
  \label{tabla:#4}
  \end{center}
}


\graphicspath{ {./img/} }

% Capítulos
\chapterstyle{bianchi}
\newcommand{\capitulo}[2]{
	\setcounter{chapter}{#1}
	\setcounter{section}{0}
	\cleardoublepage\chapter*{{\color{blue} #1:} #2}
	\addcontentsline{toc}{chapter}{#1. #2}
	\markboth{#2}{#2}
}

% Apéndices
\renewcommand{\appendixname}{Apéndice}
\renewcommand*\cftappendixname{\appendixname}

\newcommand{\apendice}[1]{
	%\renewcommand{\thechapter}{A}
	\chapter{#1}
}

\renewcommand*\cftappendixname{\appendixname\ }

% Formato de portada
\makeatletter
\newcommand{\tutor}[1]{\def\@tutor{#1}}
\newcommand{\course}[1]{\def\@course{#1}}

%\usepackage{setspace}
\def\maketitle{
  \null
  \thispagestyle{empty}
  % Cabecera ----------------
    \begin{center}%
    	{\noindent\Huge \textsc{UCJC}}\\[.3cm]%
    	{\noindent\Large \textsc{Grado en Informática}}\\[.3cm]%
    
    	\includegraphics[height=3cm]{img/escudoUCJC} \hfill
    	\includegraphics[height=3cm]{img/escudoInfor.pdf} \hfill
    \end{center}%
  \vfill
    \begin{center}%
      % Título proyecto y escudo informática ----------------
      {\noindent\large\textsc{Trabajo Fin de Grado}}\\[.3cm]
      \colorbox{ctitbox}{%
        \begin{minipage}{.95\textwidth}
          \vspace{.2cm}\large
        %   \begin{center}
        %   \textsc{Máster en Ingeniería Informática}
        %   \end{center}
          \begin{center}
%          \renewcommand{\baselinestretch}{2.0}
          \textbf{\LARGE\@title{}}
%          \renewcommand{\baselinestretch}{1.0}
          \end{center}
          \vspace{.2cm}
        \end{minipage}
      }%
    \end{center}%
  \hfill
  \vfill
  % Datos de alumno, curso y tutores ------------------
  \begin{center}%
  {%
    \noindent\large
    Realizado por \Large\textbf{\@author{}}\\[0.2cm]
    \maltese\hspace{.1cm}\maltese\hspace{.1cm}\maltese\\[0.2cm]
    {\Large\sffamily\bfseries UCJC}\\
    \textbf{\@date{}}\\
    Tutor: \@tutor{}\\
  }%
  \end{center}%
  \null
  \cleardoublepage
}
\makeatother


% Evitar solapes en el header
\nouppercaseheads

% Datos de portada
\title{TITULO PORTADA}
\newcommand{\nombre}[0]{José Gabriel Suárez Barreto} %%% cambio de comando
\newcommand{\depto}[0]{Informática}
\newcommand{\areac}[0]{}
\author{\nombre}
\tutor{TUTOR}
\date{\today}

\begin{document}

\maketitle


\newpage\null\thispagestyle{empty}\newpage


%%%%%%%%%%%%%%%%%%%%%%%%%%%%%%%%%%%%%%%%%%%%%%%%%%%%%%%%%%%%%%%%%%%%%%%%%%%%%%%%%%%%%%%%
\thispagestyle{empty}


\noindent
\begin{center}%
	{\noindent\Huge UCJC}\vspace{.5cm}%
	
\begin{center}%
	\includegraphics[height=3cm]{img/escudoUCJC} \hspace{1cm}
\end{center}%

	{\noindent\Large \textbf{Grado en Informática}}\vspace{.5cm}%
\end{center}%



\noindent D. \makeatletter\@tutor\makeatother, profesor del departamento de \depto, área de \areac.

\noindent \textbf{Expone}:

\noindent Que el alumno D. \nombre, %con DNI dni,
ha realizado el Trabajo final de Máster en Ingeniería Informática titulado "\makeatletter\textsc{\@title{}}\makeatother". 

\noindent Y que dicho trabajo ha sido realizado por el alumno bajo la dirección del que suscribe, en virtud de lo cual se autoriza su presentación y defensa.

\begin{center} %\large
En Madrid, {\large \today}
\end{center}

\vfill\vfill\vfill

% Author and supervisor
\begin{minipage}{0.45\textwidth}
\begin{flushleft} %\large
Vº. Bº. del Tutor:\\[2cm]
D. nombre tutor
\end{flushleft}
\end{minipage}
\hfill
\begin{minipage}{0.45\textwidth}
\begin{flushleft} %\large
Vº. Bº. del co-tutor:\\[2cm]
D. nombre co-tutor
\end{flushleft}
\end{minipage}
\hfill

\vfill

% para casos con solo un tutor comentar lo anterior
% y descomentar lo siguiente
%Vº. Bº. del Tutor:\\[2cm]
%D. nombre tutor


\newpage\null\thispagestyle{empty}\newpage




\frontmatter

% Abstract en castellano
\renewcommand*\abstractname{Resumen}
\begin{abstract}
En este primer apartado se hace una \textbf{breve} presentación del tema que se aborda en el proyecto.
\end{abstract}

\renewcommand*\abstractname{Descriptores}
\begin{abstract}
Palabras separadas por comas que identifiquen el contenido del proyecto Ej: servidor web, buscador de vuelos, android \ldots
\end{abstract}

\clearpage

% Abstract en inglés
\renewcommand*\abstractname{Abstract}
\begin{abstract}
A \textbf{brief} presentation of the topic addressed in the project.
\end{abstract}

\renewcommand*\abstractname{Keywords}
\begin{abstract}
keywords separated by commas.
\end{abstract}

\clearpage

% Indices
\tableofcontents
\clearpage
\listoffigures
\clearpage
\listoftables
\clearpage
\mainmatter

\capitulo{1}{Introducción}

Descripción del contenido del trabajo y del estrucutra de la memoria y del resto de materiales entregados.

\capitulo{2}{Objetivos del proyecto}

Este apartado explica de forma precisa y concisa cuales son los objetivos que se persiguen con la realización del proyecto. Se puede distinguir entre los objetivos marcados por los requisitos del software a construir y los objetivos de carácter técnico que plantea a la hora de llevar a la práctica el proyecto.

\capitulo{2}{Conceptos teóricos}

El desarrollo de un chatbot que automatice el proceso de triaje médico se presenta como una solución viable debido al avance de tecnologías como la inteligencia artificial y el procesamiento de lenguaje natural. Estos avances, combinados con la necesidad de optimizar la atención primaria y reducir la saturación en los servicios de salud, respaldan la viabilidad de desarrollar e implementar un chatbot para el triaje médico.

Además, la integración de herramientas como Amazon Comprehend Medical permite extraer y analizar datos clínicos de manera segura y eficiente, garantizando el cumplimiento de las normativas de protección de datos. Asimismo, el uso de modelos avanzados de inteligencia artificial mejora la capacidad del chatbot para comprender y responder a las consultas de los usuarios con precisión y eficiencia, optimizando el proceso de clasificación y priorización de pacientes.

\section{Evolución del Triaje Médico}

El triaje es un proceso de clasificación que nos permite identificar y priorizar a los pacientes en función de la gravedad de su estado de salud. El objetivo principal es una asistencia eficaz y eficiente, y por tanto, una herramienta rápida, fácil de aplicar y que además poseen un fuerte valor predictivo de gravedad, de evolución y de utilización de recursos. \cite{SOLER2010}

El origen del triaje se remonta a la época napoleónica, cuando se implementó por primera vez como un mecanismo para clasificar a los heridos en el campo de batalla según la urgencia de su atención. No fue hasta 1964 cuando este sistema se adoptó en el ámbito civil, estableciendo una clasificación inicial de tres niveles: pacientes no urgentes, urgentes y emergentes. A partir de la década de 1990, el triaje evolucionó a una escala de cinco niveles, expandiéndose a diversos países con sistemas propios, como la Australia Triage Scale en Australia, el Canadian Emergency Department Triage and Acuity Scale en Canadá y el Manchester Triage System en el Reino Unido, desarrollado en la ciudad homónima.

En España, el Sistema Estructurado de Triaje (SET) se ha consolidado como el estándar en la clasificación de pacientes en los servicios de urgencias. Este será el sistema base sobre el cual operará el chatbot desarrollado en este proyecto. \cite{eswiki_triaje}

Además de la clasificación por niveles de urgencia, el triaje emplea herramientas complementarias para evaluar la gravedad de los pacientes. Entre ellas, destacan las escalas de evaluación del dolor, siendo la Escala Numérica de Intensidad del Dolor (ENV) una de las más utilizadas. En este proyecto, se implementará esta escala para mejorar la precisión en la valoración del estado del paciente, permitiendo una clasificación más efectiva en función de la intensidad del dolor percibido.
Este método consiste en que el paciente puntúe su dolor del 0 al 10, siendo 0 la ausencia de dolor y 10 el peor dolor imaginable. \cite{EscalaDolorCM}
\subsection{Sistema Estructurado de Triaje}

El \textbf{Sistema Estructurado de Triaje} es un modelo de clasificación que organiza a los pacientes en cinco categorías según la gravedad de su cuadro clínico. Su principal objetivo es priorizar la urgencia del paciente sobre cualquier otro criterio estructural o profesional, garantizando una atención eficiente dentro del modelo de especialización en urgencias. Este sistema evalúa tanto el nivel de urgencia como la gravedad del caso, considerando el motivo de consulta como un factor determinante en la clasificación. \cite{SET_triaje}\cite{triaje_espana}

\section{Aplicación de inteligencia artificial en el sector de la salud}
\capitulo{4}{Técnicas y herramientas}

Esta parte de la memoria tiene como objetivo presentar las técnicas metodológicas y las herramientas de desarrollo que se han utilizado para llevar a cabo el proyecto. Si se han estudiado diferentes alternativas de metodologías, herramientas, bibliotecas se puede hacer un resumen de los aspectos más destacados de cada alternativa, incluyendo comparativas entre las distintas opciones y una justificación de las elecciones realizadas. 

No se pretende que este apartado se convierta en un capítulo de un libro dedicado a cada una de las alternativas, sino comentar los aspectos más destacados de cada opción, con un repaso somero a los fundamentos esenciales y referencias bibliográficas para que el lector pueda ampliar su conocimiento sobre el tema.



\capitulo{5}{Aspectos relevantes del desarrollo del proyecto}

Este apartado pretende recoger los aspectos más interesantes del desarrollo del proyecto, comentados por los autores del mismo.
Debe incluir desde la exposición del ciclo de vida utilizado, hasta los detalles de mayor relevancia de las fases de análisis, diseño e implementación.
Se busca que no sea una mera operación de copiar y pegar diagramas y extractos del código fuente, sino que realmente se justifiquen los caminos de solución que se han tomado, especialmente aquellos que no sean triviales.
Puede ser el lugar más adecuado para documentar los aspectos más interesantes del diseño y de la implementación, con un mayor hincapié en aspectos tales como el tipo de arquitectura elegido, los índices de las tablas de la base de datos, normalización y desnormalización, distribución en ficheros3, reglas de negocio dentro de las bases de datos (EDVHV GH GDWRV DFWLYDV), aspectos de desarrollo relacionados con el WWW...
Este apartado, debe convertirse en el resumen de la experiencia práctica del proyecto, y por sí mismo justifica que la memoria se convierta en un documento útil, fuente de referencia para los autores, los tutores y futuros alumnos.

\section{Planificación del proyecto}

\capitulo{6}{Trabajos relacionados}

Este apartado sería parecido a un estado del arte de una tesis o tesina. En un trabajo final de máster no parece tan obligada su presencia, aunque se puede dejar a juicio del tutor el incluir un pequeño resumen comentado de los trabajos y proyectos ya realizados en el campo del proyecto en curso. 

\capitulo{7}{Conclusiones y Líneas de trabajo futuras}

Todo proyecto debe incluir las conclusiones que se derivan de su desarrollo. Éstas pueden ser de diferente índole, dependiendo de la tipología del proyecto, pero normalmente van a estar presentes un conjunto de conclusiones relacionadas con los resultados del proyecto y un conjunto de conclusiones técnicas. 
Además, resulta muy útil realizar un informe crítico indicando cómo se puede mejorar el proyecto, o cómo se puede continuar trabajando en la línea del proyecto realizado. 



%\renewcommand\chaptername{Anexo}
%\renewcommand\thechapter{\Roman{chapter}}
%\setcounter{chapter}{0}

% Añadir entrada en el índice: Anexos
\appendix
\addcontentsline{toc}{part}{Apéndices}
\part*{Apéndices}

\apendice{Plan de Proyecto}

Este apéndice presentará el plan de proyecto elaborado para la realización del trabajo. En el caso de trabajos que supongan el desarrollo de software, será sustituído por el Plan de Desarrollo de Software.

\section{Introducción}

\section{Planificación temporal}

\section{Estudio de viabilidad}

\subsection{Viabilidad económica}

\subsection{Viabilidad legal}



\apendice{Especificación de Requisitos}

Si el TFM comporta el desarrollo de software, en este apéndice se reunirá la Especificación de Requisitos del mismo.

\section{Introducción}

\section{Objetivos generales}

\section{Catalogo de requisitos}

\section{Especificación de requisitos}



\apendice{Documento de Diseño}

Si el TFM comporta el desarrollo de software, en este apéndice se reunirá el Documento de Diseño asociado al mismo.

\section{Introducción}

\section{Diseño de datos}

\section{Diseño procedimental}

\section{Diseño arquitectónico}



\apendice{Documentación del Programador}

En este apéndice se incluye la documentación necesaria para que el programador de aplicaciones pueda comprender la estructura de la solución software aportada y para poder modificarla.


\section{Introducción}

\section{Estructura de directorios}

\section{Manual del programador}

\section{Compilación, instalación y ejecución del proyecto}

\section{Pruebas del sistema}

\apendice{Documentación de usuario}

En este apéndice se incluye la documentación necesaria para que el usuario sepa cómo debe instalar y usar el sistema desarrollado.

\section{Introducción}

\section{Requisitos de usuarios}

\section{Instalación}

\section{Manual del usuario}





\bibliographystyle{acm}
\bibliography{bibliografia}

\end{document}
